\section{PENDAHULUAN}
Judul adalah ringkasan berita dan menjadi dasar pertama bagi pembaca untuk memutuskan membaca isi dan memercayai pemberitaan. Interaksi pembaca daring kini bertumpu pada pemindaian cepat sehingga judul dibaca tanpa isi dibaca sampai tuntas. Ketika judul tidak mencerminkan peristiwa yang diberitakan, pembaca memperoleh gambaran yang salah. Risiko salah paham membesar ketika rentang perhatian menyusut dan semakin membesar ketika kebiasaan verifikasi masih lemah. Pada isu sensitif seperti SARA, salah paham yang kecil terhadap suatu berita yang demikian dapat berubah menjadi kegaduhan di masyarakat.

Persoalan ini terlihat sepele tetapi sulit diselesaikan secara otomatis. Di dunia akademik, permasalahan ini lebih dikenal sebagai \textit{Headline Incongruity Detection} (\textbf{\texttt{HID}}). Metode paling sederhana dari HID adalah dengan memeriksa apakah kata pada judul muncul di dalam isi. Akan tetapi, cara tersebut gagal ketika judul dan isi memakai kata yang sama untuk peristiwa yang bertolak belakang \cite{chesney17}. Penghitungan skor kemiripan kata juga gagal pada parafrasa dan pada perbedaan kecil nama atau angka yang mengubah makna \cite{yoon21, bourgonje17}. Kegagalan serupa pun juga terdokumentasi pada data berita yang berbahasa Indonesia \cite{william20, surianto26}. Penentuan kesesuaian antara judul berita dan teks isi menuntut analisis semantik yang mendalam. Penilaian tidak bisa hanya berdasarkan frekuensi kata, melainkan memerlukan analisis konteks dari judul dengan isi berita. Maka dari itu, diperlukan Sistem \textbf{\texttt{HID}} yang \textit{robust} terhadap keragaman teks yang tersebar pada media digital yang dapat menangkap makna semantik dari judul berita dengan isinya, lalu membandingkannya sebagai dasar keputusan untuk dapat menentukan kesesuaian judul berita dengan isinya.

\textbf{Latar Belakang.} Tantangan HID pada penelitian ini dapat dirangkum ke dalam beberapa poin: (1) Distribusi label sangat tidak seimbang. Data latih memuat 14397 pasangan dengan 12960 berlabel \texttt{Sesuai} dan 1437 berlabel \texttt{Tidak Sesuai}. (2) Panjang judul dan isi timpang. Judul rata-rata 10 kata. Isi rata-rata 310 kata dengan terdapat pencilan yang memiliki kata sangat panjang. (3) Kemiripan kata tinggi di kedua kelas. (4) Topik Covid yang mendominasi judul pada dataset dan aturan lomba yang mewajibkan model harus dilatih \textit{from scratch} tanpa bobot pra latih.

Pendekatan awal mengandalkan fitur artifisial dan TF-IDF. Riedel \textit{et al.} membangun baseline FNC-1 yang sulit dikalahkan dengan TF-IDF dan SVD serta XGBoost \cite{riedel17}. Yoon \textit{et al.} kemudian mengembangkan encoder hierarkis ganda berbasis perhatian untuk memodelkan paragraf isi yang panjang \cite{yoon19}. Selain itu, Mishra \textit{et al.} mengusulkan pencocokan semantik dengan perhatian timbal balik antara judul asli dan judul sintetik \cite{mishra20}S Sayangnya, seluruh pendekatan tersebut menunjukkan degradasi pada kondisi dataset ini. Baseline leksikal gagal membedakan tukar nama dan angka yang mempertahankan kata sama, sedangkan di satu sisi, model hierarkis dan semantik membutuhkan data besar dan representasi pra latih sehingga tidak dapat dipakai karena \textit{constraint} dari penelitian ini.

\textbf{Tujuan, Manfaat, dan Kontribusi.} Penelitian ini bertujuan membangun sistem HID yang robust terhadap variasi panjang teks dan kemiripan judul isi yang tinggi melalui ensemble GBDT, dilengkapi dengan fitur penyelarasan dan perbandingan rival yang dikombinasikan dengan cross encoder BiGRU. Kontribusi teknis penelitian ini meliputi tiga hal berikut.
\begin{itemize}
    \item Monotonic alignment (ALG) dan aligned window (ALI) yang menangkap pertukaran peran dan urutan kata antara judul dan isi.
    \item Paired rival comparison (DUP dan NDP) yang mengubah pertanyaan menjadi judul mana yang lebih berhak atas suatu isi. Fitur ini didukung rival ranking (RET) berbasis indeks TF-IDF.
    \item Cross encoder BiGRU dengan synthetic corruption generator tujuh resep yang melatih model peka terhadap permutasi dan negasi yang lolos dari handcrafted features.
\end{itemize}
Terdapat dua manfaat dari penelitian ini. Pertama, bagi komunitas riset dan redaksi, sistem ini diharapkan menjadi acuan terbuka pengembangan deteksi Headline Incongruity berbahasa Indonesia sekaligus alat bantu kurasi judul sebelum diterbitkan. Kedua, bagi masyarakat, sistem ini menjadi pelindung dari salah paham akibat judul yang menyesatkan.

\textbf{Batasan Masalah.} Penelitian ini dibatasi pada empat ruang lingkup berikut. Pertama, evaluasi dilakukan secara eksklusif pada dataset kompetisi DAC IFEST 2026 tanpa pengujian lintas dataset. Kedua, tugas dibatasi pada klasifikasi biner Sesuai dan Tidak Sesuai untuk berita berbahasa Indonesia dan tidak mencakup jenis misinformasi lain di luar ketidaksesuaian judul dan isi. Ketiga, seluruh parameter diinisialisasi acak dan dioptimasi pada korpus latih tanpa memanfaatkan bobot pra latih apapun dan data uji hanya dipakai pada tahap inferensi. Keempat, sistem dioptimasi untuk Macro F1 kompetisi dan belum dioptimasi untuk inferensi real-time.



